\documentclass[a4paper,11pt]{article}

% --- Packages ---
\usepackage[utf8]{inputenc}
\usepackage[T1]{fontenc}
\usepackage[french]{babel}
\usepackage{geometry}
\geometry{margin=2.5cm}
\usepackage{graphicx}
\usepackage{xcolor}
\usepackage{amsmath}
\usepackage{amssymb}
\usepackage{booktabs}
\usepackage{array}
\usepackage{enumitem}
\usepackage{hyperref}
\usepackage{fancyhdr}
\usepackage{titlesec}
\usepackage{tcolorbox}
\tcbuselibrary{skins, breakable}
\usepackage{siunitx}
\sisetup{locale = FR}

% --- Colors ---
\definecolor{bmsprimary}{HTML}{1B4F72}
\definecolor{bmssecondary}{HTML}{2E86C1}
\definecolor{bmsgreen}{HTML}{1E8449}
\definecolor{bmsorange}{HTML}{E67E22}
\definecolor{bmsgray}{HTML}{7F8C8D}
\definecolor{bmslight}{HTML}{EBF5FB}
\definecolor{bmsred}{HTML}{C0392B}

% --- Header/Footer ---
\pagestyle{fancy}
\fancyhf{}
\fancyhead[L]{\small\textcolor{bmsprimary}{\textbf{Spécifications Techniques -- BMS Intelligent}}}
\fancyhead[R]{\small\textcolor{bmsgray}{Samsung SDI 60Ah}}
\fancyfoot[C]{\thepage}
\renewcommand{\headrulewidth}{0.4pt}

% --- Title formatting ---
\titleformat{\section}{\Large\bfseries\color{bmsprimary}}{}{0em}{\thesection\quad}
\titleformat{\subsection}{\large\bfseries\color{bmssecondary}}{}{0em}{\thesubsection\quad}
\titleformat{\subsubsection}{\normalsize\bfseries\color{bmsorange}}{}{0em}{}

% --- Boxes ---
\newtcolorbox{specbox}[1][]{
  colback=bmslight, colframe=bmsprimary,
  fonttitle=\bfseries, title=#1, breakable, left=5pt, right=5pt
}
\newtcolorbox{validbox}[1][]{
  colback=bmsgreen!5, colframe=bmsgreen,
  fonttitle=\bfseries, title=#1, breakable, left=5pt, right=5pt
}
\newtcolorbox{warnbox}[1][]{
  colback=bmsorange!8, colframe=bmsorange,
  fonttitle=\bfseries, title=#1, breakable
}

% ===========================================================================
\begin{document}

% --- Title ---
\begin{titlepage}
\centering
\vspace*{2cm}
{\LARGE\bfseries\textcolor{bmsprimary}{Spécifications Techniques}}\\[0.4cm]
{\Large\textcolor{bmssecondary}{Modélisation d'un BMS Intelligent}}\\[0.3cm]
{\large\textcolor{bmssecondary}{Pack Batterie 96S -- Samsung SDI 60Ah Prismatic}}\\[1.5cm]

\begin{tcolorbox}[colback=bmslight, colframe=bmsprimary, width=13cm]
\centering
\textbf{Projet}~: BMS Intelligent pour Véhicules Électriques\\[4pt]
\textbf{Cellule}~: Samsung SDI 60Ah Prismatic (NMC)\\[4pt]
\textbf{Configuration}~: 96S1P ($\approx 355$~V nominale)\\[4pt]
\textbf{Outils}~: MATLAB/Simulink R2024b, Python 3.10+\\[4pt]
\textbf{Norme}~: ISO 26262 (ASIL)
\end{tcolorbox}

\vfill
{\large\textcolor{bmsgray}{2 avril 2026}}
\end{titlepage}

\tableofcontents
\newpage

% ===========================================================================
\section{Spécifications de la Cellule}
% ===========================================================================

\begin{center}
\renewcommand{\arraystretch}{1.4}
\begin{tabular}{p{5.5cm} p{3.5cm} p{3cm}}
\toprule
\textbf{Paramètre} & \textbf{Valeur} & \textbf{Unité} \\
\midrule
Chimie & NMC & -- \\
Format & Prismatique & -- \\
Capacité nominale & 60 & Ah \\
Tension nominale & 3.7 & V \\
Tension min (cutoff) & 2.7 & V \\
Tension max (charge) & 4.2 & V \\
Courant max continu (décharge) & 60 (1C) & A \\
Courant max crête (décharge) & 180 (3C) & A \\
Plage de température opérationnelle & $-20$ à $+60$ & $^\circ$C \\
Masse cellule & $\approx 1.3$ & kg \\
\bottomrule
\end{tabular}
\end{center}

% ===========================================================================
\section{Configuration du Pack Batterie}
% ===========================================================================

\begin{center}
\renewcommand{\arraystretch}{1.4}
\begin{tabular}{p{5.5cm} p{3.5cm} p{3cm}}
\toprule
\textbf{Paramètre} & \textbf{Valeur} & \textbf{Unité} \\
\midrule
Topologie & 96S1P & -- \\
Nombre de modules & 8 & -- \\
Cellules par module & 12 (série) & -- \\
Tension nominale pack & 355.2 & V \\
Tension min pack & 259.2 & V \\
Tension max pack & 403.2 & V \\
Énergie nominale & 21.3 & kWh \\
Dispersion capacité inter-cellules & $\pm 2\%$ & -- \\
Dispersion $R_0$ inter-cellules & $\pm 5\%$ & -- \\
\bottomrule
\end{tabular}
\end{center}

\begin{specbox}[Architecture BMS]
Le BMS adopte une architecture \textbf{Master--Slave}~:
\begin{itemize}[nosep]
  \item \textbf{1~Master}~: estimation SoC/SoH, gestion des défauts, interface CAN, gestion thermique.
  \item \textbf{8~Slaves}~: un par module 12S -- mesure de tension cellulaire, équilibrage passif.
\end{itemize}
\end{specbox}

% ===========================================================================
\section{Modèle Électrique -- ECM 2RC}
% ===========================================================================

\subsection{Équation du circuit}

Le modèle de cellule est un circuit équivalent à deux boucles RC~:
\[
  V_{terminal} = OCV(SoC, T) - I \cdot R_0(SoC, T) - V_{RC1} - V_{RC2}
\]
avec les dynamiques~:
\[
  \frac{dV_{RC_i}}{dt} = \frac{I}{C_i(SoC,T)} - \frac{V_{RC_i}}{R_i(SoC,T) \cdot C_i(SoC,T)}, \quad i \in \{1, 2\}
\]

L'état de charge est mis à jour par Coulomb counting~:
\[
  SoC_{k+1} = SoC_k - \frac{I_k \cdot \Delta t}{Q_{nom} \cdot 3600}
\]
avec $Q_{nom} = 60$~Ah et $\Delta t = 0.1$~s.

\subsection{Tables de paramètres (LUT)}

\begin{center}
\renewcommand{\arraystretch}{1.4}
\begin{tabular}{lccc}
\toprule
\textbf{Paramètre} & \textbf{Grille} & \textbf{Plage typique (cellule)} & \textbf{Source} \\
\midrule
$OCV(SoC, T)$ & $101 \times 7$ & 2.86 -- 4.20~V & Mesure \\
$R_0(SoC, T)$ & $21 \times 7$ & 0.4 -- 1.6~m$\Omega$ & Mesure \\
$R_1(SoC, T)$ & $21 \times 7$ & 0.1 -- 0.4~m$\Omega$ & Mesure \\
$C_1(SoC, T)$ & $21 \times 7$ & 3\,000 -- 27\,000~F & Mesure \\
$R_2(SoC, T)$ & $21 \times 7$ & 0.4 -- 2.0~m$\Omega$ & Mesure \\
$C_2(SoC, T)$ & $21 \times 7$ & 42\,000 -- 74\,000~F & Mesure \\
\bottomrule
\end{tabular}
\end{center}

\textbf{Breakpoints~:}
\begin{itemize}[nosep]
  \item SoC~: $[0, 5, 10, \ldots, 100]\%$ -- 21 points (OCV~: 101 points, pas de 1\%).
  \item Température~: $[0, 5, 10, 15, 20, 25, 30]\,^\circ$C -- 7 points.
  \item Interpolation~: bilinéaire dans Simulink.
  \item Lissage gaussien appliqué~: $\sigma = [0.8, 0.5]$ (SoC, T).
\end{itemize}

\subsection{Méthode d'extraction}

\begin{center}
\renewcommand{\arraystretch}{1.3}
\begin{tabular}{lp{8cm}}
\toprule
\textbf{Paramètre} & \textbf{Méthode} \\
\midrule
$OCV$ & Médiane des tensions en phase quasi-statique ($|I| < 3$~A, durée $> 3$~s) par bin $(SoC, T)$. \\
$R_0$ & Saut de tension instantané lors d'échelons de courant ($|\Delta I| > 8$~A). \\
$R_1, C_1, R_2, C_2$ & Ajustement par moindres carrés sur fenêtres de 60~s (stride 30~s, RMSE $< 15$~mV). \\
\bottomrule
\end{tabular}
\end{center}

% ===========================================================================
\section{Modèle Thermique}
% ===========================================================================

\subsection{Équation thermique par cellule}

\[
  C_{th}\,\frac{dT_{cell}}{dt}
  = \underbrace{I^2 \cdot (R_0 + R_1 + R_2)}_{P_{Joule}}
  - h_{cool}\,(T_{cell} - T_{cool})
  - h_{amb}\,(T_{cell} - T_{amb})
\]

\subsection{Paramètres thermiques calibrés}

\begin{center}
\renewcommand{\arraystretch}{1.4}
\begin{tabular}{lrl}
\toprule
\textbf{Paramètre} & \textbf{Valeur} & \textbf{Description} \\
\midrule
$C_{th}$ & 1\,300~J/K & Capacité thermique cellule \\
$h_{cool}$ & 0.02~W/K & Couplage cellule--liquide de refroidissement \\
$h_{amb}$ & 0.15~W/K & Échange cellule--ambiant \\
\bottomrule
\end{tabular}
\end{center}

\subsection{Circuit de refroidissement (\texttt{Coolant\_loop})}

\[
  C_{cool}\,\frac{dT_{cool}}{dt}
  = \dot{m}_{cool}\,c_p\,(T_{cool,in} - T_{cool}) + Q_{cell \to cool}
\]
Entrée~: $T_{cool,in}(t)$ issu du signal mesuré \texttt{coolant\_temp\_inlet\_c}
(plage~: $-1$ à $70\,^\circ$C dans le dataset).

% ===========================================================================
\section{Domaine de Validité du Modèle}
% ===========================================================================

\begin{validbox}[Enveloppe opérationnelle validée]
Le modèle ECM 2RC du pack 96S est \textbf{validé dans l'enveloppe suivante},
directement issue de la couverture du dataset de mesure~:
\begin{itemize}[nosep]
  \item $SoC \in [40\%, 90\%]$~: plage effectivement parcourue et représentée
    dans les données de mesure utilisées pour l'identification des paramètres.
  \item $T_{cell} \in [10\,^\circ C, 30\,^\circ C]$~: plage à forte densité de
    points de mesure.
\end{itemize}
\end{validbox}

\subsection{Performance dans l'enveloppe}

\begin{center}
\renewcommand{\arraystretch}{1.4}
\begin{tabular}{lccc}
\toprule
\textbf{Plage SoC} & \textbf{RMSE~V typique (pack)} & \textbf{RMSE par cellule} & \textbf{Qualité} \\
\midrule
50 -- 90\% & 200 -- 700~mV & $\approx 2$--$7$~mV & Bon \\
40 -- 50\% & 800 -- 1\,600~mV & $\approx 8$--$17$~mV & Acceptable \\
$<40\%$ & 2\,500 -- 6\,000~mV & $\approx 26$--$63$~mV & Hors domaine \\
\bottomrule
\end{tabular}
\end{center}

\subsection{Métriques globales (70 trips)}

\begin{center}
\renewcommand{\arraystretch}{1.4}
\begin{tabular}{lcccc}
\toprule
\textbf{Indicateur} & \textbf{Médiane} & \textbf{Moyenne} & \textbf{Max} & \textbf{Unité} \\
\midrule
RMSE tension pack & 1\,102 & 1\,563 & 6\,201 & mV \\
Biais tension & -- & $-591$ & -- & mV \\
RMSE température & 1.18 & 1.46 & 4.59 & $^\circ$C \\
RMSE SoC & 0.18 & 0.25 & 0.65 & \% \\
\bottomrule
\end{tabular}
\end{center}

\begin{warnbox}[Conditions hors-enveloppe]
Les conditions hors-enveloppe ($SoC < 40\%$, $T < 10\,^\circ$C) correspondent
à des situations peu couvertes dans le dataset. Les données de simulation
ne couvrent que 3~températures ambiantes ($-10$, $0$, $+10\,^\circ$C) et ne
permettent pas de combler les bins manquants aux températures intermédiaires
(15--25$\,^\circ$C). L'extension du modèle à ces plages nécessiterait des
données de caractérisation dédiées (pulses HPPC à basse température,
décharges complètes contrôlées).
\end{warnbox}

% ===========================================================================
\section{Dataset Utilisé}
% ===========================================================================

\begin{center}
\renewcommand{\arraystretch}{1.3}
\begin{tabular}{lrrr}
\toprule
\textbf{Dataset} & \textbf{Fichiers} & \textbf{Lignes} & \textbf{Colonnes} \\
\midrule
Mesure -- TripA (été) & 32 & 467\,701 & 22--28 \\
Mesure -- TripB (hiver) & 38 & 627\,092 & 48 \\
Simulation (CAN + LIN) & 24 & 917\,281 & 43 \\
\midrule
\textbf{Total} & \textbf{94} & \textbf{2\,012\,074} & -- \\
\bottomrule
\end{tabular}
\end{center}

\begin{center}
\renewcommand{\arraystretch}{1.3}
\begin{tabular}{lll}
\toprule
\textbf{Signal} & \textbf{TripA (été)} & \textbf{TripB (hiver)} \\
\midrule
Tension pack & 349 -- 395~V & 302 -- 394~V \\
Courant batterie & $-395$ à $+144$~A & $-404$ à $+145$~A \\
Température cellule & 16 -- 32$\,^\circ$C & $-1$ -- 22$\,^\circ$C \\
Température ambiante & 14 -- 33.5$\,^\circ$C & $-3.5$ -- 14$\,^\circ$C \\
SoC & 34 -- 89\% & 0 -- 86\% \\
Fréquence d'échantillonnage & \multicolumn{2}{c}{10~Hz ($\Delta t = 0.1$~s)} \\
\bottomrule
\end{tabular}
\end{center}

% ===========================================================================
\section{Spécifications du Contrôleur BMS (Phase 3 -- à implémenter)}
% ===========================================================================

\begin{center}
\renewcommand{\arraystretch}{1.4}
\begin{tabular}{p{4cm} p{8.5cm}}
\toprule
\textbf{Fonction} & \textbf{Spécification cible} \\
\midrule
Estimation SoC (EKF) &
  État~: $[SoC, V_{RC1}, V_{RC2}]^T$. Objectif~: RMSE $< 2\%$. \\
Estimation SoH &
  Tracking $R_0$ par RLS. $SoH_R = R_{0,init} / R_{0,actuel}$. \\
Équilibrage passif &
  Seuil~: $\Delta V > 20$~mV. Stateflow 4 états. \\
Détection de défauts &
  Surtension ($> 4.25$~V), sous-tension ($< 2.7$~V),
  surcourant ($> 180$~A), surchauffe ($> 55\,^\circ$C),
  sous-température ($< -20\,^\circ$C), défaut capteur. Debounce~: $N$ échantillons. \\
Gestion thermique &
  PI~: chauffage si $T < 5\,^\circ$C, refroidissement si $T > 40\,^\circ$C. \\
Interface CAN &
  Messages~: BMS\_Status, BMS\_Faults, BMS\_Limits, BMS\_Thermal. \\
\bottomrule
\end{tabular}
\end{center}

% ===========================================================================
\section{Spécifications du Module IA (Phase 4 -- à implémenter)}
% ===========================================================================

\begin{center}
\renewcommand{\arraystretch}{1.4}
\begin{tabular}{p{4cm} p{8.5cm}}
\toprule
\textbf{Élément} & \textbf{Spécification} \\
\midrule
Architecture &
  LSTM 128 unités + Dense 64/32 + Softmax 6 classes. \\
Entrée &
  Fenêtre 100 pas $\times$ 4 features ($V, I, T, SoC$). \\
Classes &
  Overvoltage, Undervoltage, Overcurrent, Overtemperature,
  SensorFault, InternalShortCircuit. \\
Objectif &
  F1-score $> 0.95$ sur set de test. \\
Export &
  ONNX vers Simulink (Deep Learning Toolbox). \\
\bottomrule
\end{tabular}
\end{center}

% ===========================================================================
\section{Spécifications de Validation (Phase 5 -- à implémenter)}
% ===========================================================================

\begin{center}
\renewcommand{\arraystretch}{1.4}
\begin{tabular}{p{3cm} p{9.5cm}}
\toprule
\textbf{Niveau} & \textbf{Critères} \\
\midrule
MIL &
  70 trips + 12 scénarios simulation + défauts injectés.
  Couverture MC/DC 100\% Stateflow. \\
SIL &
  Code C (Embedded Coder, \texttt{ert.tlc}, ARM Cortex-M4).
  Écart MIL/SIL $< 10^{-6}$. Conformité MISRA~C. \\
PIL &
  MCU cible (STM32 / NXP S32K). $T_{exec} < T_s = 10$~ms.
  Mémoire $< 80\%$ capacité MCU. \\
\bottomrule
\end{tabular}
\end{center}

% ===========================================================================
\section{Arborescence des Fichiers du Projet}
% ===========================================================================

\begin{center}
\small
\renewcommand{\arraystretch}{1.2}
\begin{tabular}{p{6.5cm} p{6cm}}
\toprule
\textbf{Chemin} & \textbf{Description} \\
\midrule
\texttt{plant\_model/models/cell\_ecm\_2rc.slx} & Modèle cellule ECM 2RC \\
\texttt{plant\_model/models/module\_12s.slx} & Module 12 cellules série \\
\texttt{plant\_model/models/pack\_96s.slx} & Pack complet 96S \\
\texttt{plant\_model/data/*\_lut.csv} & Tables LUT ($R_0$--$C_2$, OCV) \\
\texttt{plant\_model/data/thermal\_params.csv} & Paramètres thermiques \\
\texttt{plant\_model/data/validation\_results.csv} & Résultats validation 70 trips \\
\texttt{plant\_model/scripts/extract\_ecm\_params.py} & Extraction paramètres ECM \\
\texttt{plant\_model/scripts/extract\_ocv.py} & Extraction courbe OCV \\
\texttt{plant\_model/scripts/calibrate\_thermal.m} & Calibration thermique \\
\texttt{plant\_model/scripts/validate\_pack\_model.m} & Validation pack \\
\texttt{ai/data\_preparation/*.py} & Nettoyage et qualité données \\
\texttt{ai/reports/*.md} & Rapports EDA, qualité, nettoyage \\
\texttt{data/cleaned/} & Datasets nettoyés \\
\bottomrule
\end{tabular}
\end{center}

\end{document}
